\documentclass[12pt,a4paper]{article}

\usepackage{cmap}
\usepackage[T2A]{fontenc}
\usepackage[utf8]{inputenc}
\usepackage[russian]{babel}
\usepackage{amsmath,amssymb}
\usepackage[a4paper,margin=2cm]{geometry}

\newcommand{\R}{\mathbb{R}}
\newcommand{\eps}{\varepsilon}
\newcommand{\Ball}[2]{U_{#1}(#2)}
\newcommand{\PBall}[2]{\mathring U_{#1}(#2)}
\newcommand{\limset}[2]{\lim_{\substack{#1\\ #2}}}

\setlength{\parindent}{0pt}
\setlength{\parskip}{0.45em}
\sloppy

\begin{document}

\begin{center}
{\Large\bfseries Билет 2}
\end{center}

\noindent\fbox{%
  \begin{minipage}{\dimexpr\linewidth-2\fboxsep-2\fboxrule\relax}
  \normalfont
Предел числовой функции нескольких переменных. Предел функции по множеству.
Пределы по направлениям. Непрерывность функции нескольких переменных в точке и
по множеству. Непрерывность сложной функции. Свойства функций, непрерывных на
компакте: ограниченность, достижимость точных нижней и верхней граней,
равномерная непрерывность. Теорема о промежуточных значениях функции,
непрерывной в области.
  \end{minipage}%
}

\section*{Краткий план}

Сначала вводится предел \(f(x)\) при \(x\to x_0\) по совокупности переменных:
через \(\eps\)-\(\delta\) и через последовательности Гейне. Затем отдельно
рассматриваются предел по множеству и предел по направлению. После этого
непрерывность в точке определяется как совпадение предела со значением функции,
а непрерывность на множестве -- поточечно. Основные теоремы: композиция
непрерывных функций непрерывна; непрерывная функция на компакте ограничена,
достигает минимума и максимума и равномерно непрерывна; непрерывная на области
скалярная функция принимает все промежуточные значения между двумя своими
значениями.

\section*{Определения}

\textbf{Числовая функция нескольких переменных.}
Пусть \(X\subset\R^n\). Функция \(f:X\to\R\) сопоставляет каждой точке
\[
  x=(x_1,\dots,x_n)\in X
\]
единственное число \(f(x)=f(x_1,\dots,x_n)\).

\textbf{Проколотая окрестность.}
Для \(x_0\in\R^n\)
\[
  \PBall{\delta}{x_0}
  =\{x\in\R^n:\ 0<|x-x_0|<\delta\}.
\]

\textbf{Предел по совокупности переменных.}
Пусть \(f\) определена в некоторой проколотой окрестности точки \(x_0\).
Число \(A\in\R\) называется пределом \(f(x)\) при \(x\to x_0\), если
\[
  \forall\eps>0\ \exists\delta>0\ \forall x\in\PBall{\delta}{x_0}:
  |f(x)-A|<\eps .
\]
Пишут
\[
  \lim_{x\to x_0}f(x)=A
  \quad\text{или}\quad
  \lim_{\substack{x_1\to x_1^0\\ \dots\\ x_n\to x_n^0}}
  f(x_1,\dots,x_n)=A.
\]
Эквивалентно, по Гейне: для любой последовательности \(x_k\ne x_0\),
\(x_k\to x_0\), выполняется \(f(x_k)\to A\).

\textbf{Предельная и изолированная точки множества.}
Точка \(x_0\) называется предельной точкой множества \(X\subset\R^n\), если
существует последовательность \(x_k\in X\), \(x_k\ne x_0\), такая что
\(x_k\to x_0\). Точка \(x_0\in X\) называется изолированной точкой \(X\), если
\[
  \exists\delta>0:\ \PBall{\delta}{x_0}\cap X=\varnothing.
\]

\textbf{Предел по множеству.}
Пусть \(x_0\) -- предельная точка \(X\subset\R^n\), \(f:X\to\R\). Число
\(A\) называется пределом \(f\) в точке \(x_0\) по множеству \(X\), если
\[
  \forall\eps>0\ \exists\delta>0\ \forall x\in X\cap\PBall{\delta}{x_0}:
  |f(x)-A|<\eps .
\]
Пишут
\[
  \limset{x\to x_0}{x\in X} f(x)=A.
\]
Эквивалентно: для любой последовательности Гейне \(x_k\in X\),
\(x_k\to x_0\), \(x_k\ne x_0\), имеем \(f(x_k)\to A\).

\textbf{Предел по направлению.}
Направление в \(\R^n\) -- это вектор \(\ell\in\R^n\) с \(|\ell|=1\). Число
\(A\) называется пределом \(f\) в точке \(x_0\) по направлению \(\ell\), если
\[
  \lim_{t\to +0} f(x_0+t\ell)=A,
\]
то есть
\[
  \forall\eps>0\ \exists\delta>0\ \forall t\in(0,\delta):
  |f(x_0+t\ell)-A|<\eps .
\]

\textbf{Непрерывность в точке по множеству.}
Функция \(f:X\to\R\) называется непрерывной в точке \(x_0\in X\) по множеству
\(X\), если либо \(x_0\) -- предельная точка \(X\) и
\[
  \limset{x\to x_0}{x\in X} f(x)=f(x_0),
\]
либо \(x_0\) -- изолированная точка \(X\).

\textbf{Непрерывность в точке по совокупности переменных.}
Если \(x_0\) -- внутренняя точка \(X\), то \(f\) непрерывна в точке \(x_0\),
если
\[
  \lim_{x\to x_0}f(x)=f(x_0),
\]
то есть
\[
  \forall\eps>0\ \exists\delta>0\ \forall x\in X:
  |x-x_0|<\delta \Rightarrow |f(x)-f(x_0)|<\eps .
\]

\textbf{Непрерывность на множестве.}
Функция \(f:X\to\R\) непрерывна на множестве \(X\), если она непрерывна в
каждой точке \(x_0\in X\) по множеству \(X\).

\textbf{Равномерная непрерывность.}
Функция \(f:X\to\R\) равномерно непрерывна на \(X\), если
\[
  \forall\eps>0\ \exists\delta>0\ \forall x,x'\in X:
  |x-x'|<\delta \Rightarrow |f(x)-f(x')|<\eps .
\]
В отличие от обычной непрерывности, одно и то же \(\delta\) работает сразу для
всех точек множества.

\textbf{Область.}
Множество \(D\subset\R^n\) называется областью, если оно открыто и линейно
связно. Линейная связность означает, что любые \(x_1,x_2\in D\) можно соединить
кривой в \(D\): существует непрерывная \(r:[a,b]\to D\), для которой
\(r(a)=x_1\), \(r(b)=x_2\).

\section*{Теоремы}

\textbf{1. Критерий Гейне для предела.}
Предел \(\lim\limits_{x\to x_0}f(x)=A\) существует тогда и только тогда, когда
для всякой последовательности \(x_k\ne x_0\), \(x_k\to x_0\), выполняется
\(f(x_k)\to A\). Такой же критерий справедлив для предела по множеству.

\textbf{2. Связь предела и пределов по направлениям.}
Если существует предел \(\lim\limits_{x\to x_0}f(x)=A\), то по любому
направлению \(\ell\) существует предел
\[
  \lim_{t\to+0}f(x_0+t\ell)=A.
\]
Обратное неверно: существование одинаковых пределов по всем направлениям само
по себе не гарантирует существование предела по совокупности переменных.

\textbf{3. Непрерывность сложной функции.}
Пусть \(X\subset\R^n\), \(Y\subset\R^m\), \(f:X\to Y\), \(g:Y\to\R^p\).
Если \(f\) непрерывна на \(X\), а \(g\) непрерывна на \(Y\), то
\[
  \varphi(x)=g(f(x))
\]
непрерывна на \(X\).

\textbf{4. Непрерывный образ компакта.}
Если \(f:X\to\R^m\) непрерывна на компакте \(X\subset\R^n\), то \(f(X)\) --
компакт.

\textbf{5. Ограниченность на компакте.}
Если \(f:X\to\R^m\) непрерывна на компакте \(X\subset\R^n\), то \(f\)
ограничена на \(X\).

\textbf{6. Теорема Вейерштрасса.}
Если скалярная функция \(f:X\to\R\) непрерывна на компакте \(X\subset\R^n\), то
существуют
\[
  \min_{x\in X}f(x), \qquad \max_{x\in X}f(x).
\]
Иными словами, точные нижняя и верхняя грани достигаются.

\textbf{7. Теорема Кантора.}
Если \(f:X\to\R\) непрерывна на компакте \(X\subset\R^n\), то \(f\) равномерно
непрерывна на \(X\).

\textbf{8. Теорема о промежуточных значениях.}
Пусть \(D\subset\R^n\) -- область, \(f:D\to\R\) непрерывна на \(D\), и
\(f(x_1)=y_1\), \(f(x_2)=y_2\). Тогда для всякого числа \(y_0\), лежащего между
\(y_1\) и \(y_2\), существует точка \(x_0\in D\), такая что
\[
  f(x_0)=y_0.
\]
Более общая формулировка верна для любого линейно связного множества \(X\).

\section*{Доказательства}

\textbf{Критерий Гейне для предела.}
Пусть выполнено определение Коши и \(x_k\to x_0\), \(x_k\ne x_0\). Для данного
\(\eps>0\) берем соответствующее \(\delta>0\). Начиная с некоторого номера
\(k\), имеем \(|x_k-x_0|<\delta\), поэтому \(|f(x_k)-A|<\eps\). Значит
\(f(x_k)\to A\).

Обратно, пусть условие Коши неверно. Тогда
\[
  \exists\eps_0>0\ \forall\delta>0\ \exists x\in\PBall{\delta}{x_0}:
  |f(x)-A|\ge\eps_0.
\]
Беря \(\delta=1/k\), строим последовательность \(x_k\ne x_0\), \(x_k\to x_0\),
но \(|f(x_k)-A|\ge\eps_0\). Это противоречит условию Гейне. Для предела по
множеству доказательство такое же, только все \(x_k\) выбираются из \(X\).

\textbf{Предел по совокупности переменных и пределы по направлениям.}
Пусть \(\lim_{x\to x_0}f(x)=A\), \(|\ell|=1\). По определению предела для
\(\eps>0\) найдется \(\delta>0\) такое, что из \(0<|x-x_0|<\delta\) следует
\(|f(x)-A|<\eps\). Если \(t\in(0,\delta)\) и \(x=x_0+t\ell\), то
\[
  |x-x_0|=t|\ell|=t<\delta,
\]
значит \(|f(x_0+t\ell)-A|<\eps\). Следовательно,
\(\lim\limits_{t\to+0}f(x_0+t\ell)=A\).

Обратное неверно. Для
\[
  f(u,v)=
  \begin{cases}
    \dfrac{uv^2}{u^2+v^4}, & (u,v)\ne(0,0),\\[0.8em]
    0, & (u,v)=(0,0),
  \end{cases}
\]
в точке \((0,0)\) предел по любой прямой \(u=t\cos\varphi\),
\(v=t\sin\varphi\) равен \(0\). Но вдоль кривой \(u=v^2\) получаем
\(f(v^2,v)=1/2\), а вдоль \(u=0\) получаем \(f(0,v)=0\). По критерию Гейне
предела по совокупности переменных нет.

\textbf{Непрерывность сложной функции.}
Зафиксируем \(x_0\in X\) и \(\eps>0\). Поскольку \(g\) непрерывна в
\(y_0=f(x_0)\), существует \(\sigma>0\) такое, что для \(y\in Y\)
\[
  |y-y_0|<\sigma \Rightarrow |g(y)-g(y_0)|<\eps.
\]
Поскольку \(f\) непрерывна в \(x_0\), существует \(\delta>0\) такое, что для
\(x\in X\)
\[
  |x-x_0|<\delta \Rightarrow |f(x)-f(x_0)|<\sigma.
\]
Следовательно,
\[
  |x-x_0|<\delta
  \Rightarrow
  |g(f(x))-g(f(x_0))|<\eps,
\]
то есть \(g\circ f\) непрерывна в \(x_0\). Так как \(x_0\) произвольна,
композиция непрерывна на \(X\).

\textbf{Непрерывный образ компакта.}
Пусть \(\{y_k\}\subset f(X)\). Для каждого \(k\) выберем \(x_k\in X\) так, что
\(f(x_k)=y_k\). Так как \(X\) компактен, из \(\{x_k\}\) можно выделить
подпоследовательность \(x_{k_j}\to x_0\in X\). По непрерывности
\[
  y_{k_j}=f(x_{k_j})\to f(x_0)\in f(X).
\]
Значит из любой последовательности в \(f(X)\) выделяется подпоследовательность,
сходящаяся к элементу \(f(X)\), то есть \(f(X)\) -- компакт.

\textbf{Ограниченность и достижение точных граней.}
Если \(f:X\to\R^m\) непрерывна на компакте, то \(f(X)\) компакт в \(\R^m\).
Всякий компакт в \(\R^m\) ограничен, поэтому \(f\) ограничена.

Пусть теперь \(f:X\to\R\). Обозначим
\[
  m=\inf_{x\in X}f(x), \qquad M=\sup_{x\in X}f(x).
\]
Так как \(f\) ограничена, \(m,M\in\R\). По определению инфимума существует
минимизирующая последовательность \(x_k\in X\), для которой \(f(x_k)\to m\).
Из компактности \(X\) выделяем \(x_{k_j}\to x_*\in X\). По непрерывности
\[
  f(x_{k_j})\to f(x_*).
\]
Следовательно, \(f(x_*)=m\), то есть точная нижняя грань достигается. Для
точной верхней грани аналогично берется максимизирующая последовательность, и
получается точка \(x^*\in X\), такая что \(f(x^*)=M\).

\textbf{Теорема Кантора.}
Пусть \(f\) непрерывна на компакте \(X\). Зафиксируем \(\eps>0\). Для каждой
точки \(a\in X\) по непрерывности найдется \(\delta(a)>0\) такое, что
\[
  x\in X,\quad |x-a|<\delta(a)
  \Rightarrow |f(x)-f(a)|<\eps/2.
\]
Семейство шаров \(\Ball{\delta(a)/2}{a}\), \(a\in X\), покрывает компакт
\(X\). Поэтому из него можно выделить конечное подпокрытие
\[
  X\subset\bigcup_{j=1}^N \Ball{\delta(a_j)/2}{a_j}.
\]
Положим
\[
  \delta=\min_{1\le j\le N}\frac{\delta(a_j)}{2}>0.
\]
Если \(x,x'\in X\) и \(|x-x'|<\delta\), то \(x\in\Ball{\delta(a_j)/2}{a_j}\)
для некоторого \(j\). Тогда
\[
  |x'-a_j|\le |x'-x|+|x-a_j|<\delta+\frac{\delta(a_j)}2\le\delta(a_j).
\]
Следовательно,
\[
  |f(x)-f(x')|
  \le |f(x)-f(a_j)|+|f(a_j)-f(x')|
  <\eps/2+\eps/2=\eps.
\]
Это и есть равномерная непрерывность.

\textbf{Теорема о промежуточных значениях в области.}
Пусть \(x_1,x_2\in D\), \(f(x_1)=y_1\), \(f(x_2)=y_2\). Так как область
линейно связна, существует непрерывная кривая \(r:[a,b]\to D\), такая что
\(r(a)=x_1\), \(r(b)=x_2\). По теореме о непрерывности сложной функции
\[
  \phi(t)=f(r(t))
\]
непрерывна на отрезке \([a,b]\). При этом \(\phi(a)=y_1\), \(\phi(b)=y_2\).
По одномерной теореме Больцано--Коши для любого \(y_0\) между \(y_1\) и
\(y_2\) найдется \(t_0\in[a,b]\), такое что \(\phi(t_0)=y_0\). Тогда
\(x_0=r(t_0)\in D\) и \(f(x_0)=y_0\).

\section*{Финальная устная версия}

Предел функции нескольких переменных \(f:X\to\R\) при \(x\to x_0\) определяется
так же, как в одномерном случае, только расстояние берется в \(\R^n\):
\[
  \forall\eps>0\ \exists\delta>0:\ 0<|x-x_0|<\delta
  \Rightarrow |f(x)-A|<\eps.
\]
Эквивалентно, по Гейне, для любой последовательности \(x_k\ne x_0\),
\(x_k\to x_0\), должно быть \(f(x_k)\to A\). Предел по множеству \(X\)
отличается только тем, что \(x\) и последовательности берутся из \(X\).

Предел по направлению \(\ell\), \(|\ell|=1\), есть одномерный предел
\(\lim_{t\to+0}f(x_0+t\ell)\). Если существует общий предел, то все пределы по
направлениям существуют и равны ему. Обратное неверно: у функции
\(\frac{uv^2}{u^2+v^4}\) в нуле все пределы по прямым равны нулю, но вдоль
\(u=v^2\) получается \(1/2\), поэтому общего предела нет.

Непрерывность в точке означает, что предел равен значению функции:
\[
  \lim_{x\to x_0}f(x)=f(x_0).
\]
По множеству \(X\) используется предел при \(x\to x_0\), \(x\in X\); если
\(x_0\) изолирована, функция считается непрерывной в этой точке. Непрерывность
на множестве -- непрерывность в каждой его точке. Композиция непрерывных
функций непрерывна: сначала малость \(f(x)-f(x_0)\) обеспечивает попадание в
окрестность \(f(x_0)\), затем непрерывность внешней функции дает малость
\(g(f(x))-g(f(x_0))\).

Если \(f\) непрерывна на компакте \(X\), то \(f(X)\) -- компакт: из
последовательности значений \(f(x_k)\) выделяем подпоследовательность
\(x_{k_j}\to x_0\in X\), и по непрерывности \(f(x_{k_j})\to f(x_0)\). Отсюда
следует ограниченность. Для скалярной функции, кроме того, достигаются
\(\inf f\) и \(\sup f\): берется минимизирующая или максимизирующая
последовательность, из нее выделяется сходящаяся подпоследовательность, и
непрерывность переносит предел на значение функции.

По теореме Кантора непрерывная на компакте функция равномерно непрерывна:
локальные \(\delta(a)\) из непрерывности дают покрытие компакта шарами,
выбирается конечное подпокрытие, и минимум конечного набора радиусов дает одно
\(\delta\) для всех точек.

Наконец, если \(D\) -- область и \(f:D\to\R\) непрерывна, то между любыми
двумя значениями \(y_1=f(x_1)\) и \(y_2=f(x_2)\) функция принимает все
промежуточные значения. Действительно, в области точки \(x_1,x_2\) соединяются
кривой \(r\), а одномерная функция \(f(r(t))\) непрерывна на отрезке, поэтому
работает теорема Больцано--Коши.

\section*{Источники}

Иванов Г.Е. \emph{Лекции по математическому анализу. Часть 1}:
\texttt{IvGE\_1.pdf}, стр. 57--60, 167--179.

Иванов Г.Е. \emph{Лекции по математическому анализу. Часть 2}:
\texttt{IvGE\_2.pdf}, стр. 57.

\end{document}
